\clearpage
\begin{explanationbox}

\begin{minipage}{\linewidth}
\textbf{\textcolor{CExplanation}{一、为什么会出现这个方法？}}
\begin{center}\begin{tikzpicture}[font=\small,>=stealth,line join=round]
\draw[gray!45,rounded corners=3pt] (0,0) rectangle (14.2,9.4);
\node at (7.1,8.8) {同样趋于零，比值仍可能是不同的数};
\begin{scope}[shift={(3.8,4.7)},scale=1.3]
\draw[->,gray] (-1.3,0)--(1.45,0) node[above] {$t=x-1$};
\draw[->,gray] (0,-2.7)--(0,2.9) node[left] {$y$};
\draw[blue!75!black,thick,samples=70,domain=-.8:.8] plot (\x,{2*\x+\x*\x});
\draw[orange!85!black,thick,dashed] (-.8,-.8)--(.8,.8);
\draw[violet,thick,samples=70,domain=-.8:.8] plot (\x,{\x*\x});
\draw[green!45!black,thick,dashdotted] (-.8,-2.4)--(.8,2.4);
\fill (0,0) circle(2pt);
\end{scope}
\node[anchor=west,align=left] at (8.4,5.2)
{$\textcolor{orange!85!black}{\text{分母 }g=t}$\\[15pt]
$\textcolor{blue!75!black}{(2t+t^2)/t\to2}$\\[15pt]
$\textcolor{violet}{t^2/t\to0}$\\[15pt]
$\textcolor{green!45!black}{3t/t\to3}$};
\node[text width=13.2cm,align=center] at (7.1,.65) {四条曲线都经过原点；“都趋于零”没有决定它们高度之比。};
\end{tikzpicture}
\end{center}
\end{minipage}



图中令 $t=x-1$：四条曲线都经过原点，但蓝、紫、绿线与橙线的比值分别趋于 $2,0,3$。
因此“分子、分母都趋于零”仍不能决定答案，必须比较它们的变化比例。

\medskip
\textbf{\textcolor{CExplanation}{二、为什么想到求导？}}

导数能描述函数在某个位置附近的变化快慢。
当 $f(a)=g(a)=0$ 时，如果两个函数在 $a$ 处都能求导，而且 $g'(a)\ne0$，那么
\[
\frac{f(x)}{g(x)}
=\frac{\frac{f(x)-f(a)}{x-a}}{\frac{g(x)-g(a)}{x-a}}
\quad(x\ne a).
\]
当 $x$ 很接近 $a$ 时，上下两个“小分式”分别接近 $f'(a)$、$g'(a)$。
于是原来的比值接近 $\frac{f'(a)}{g'(a)}$。

\textbf{眼前一亮的地方：}原本是在比两个函数的\emph{数值}，转个角度后，变成比它们的\emph{变化快慢}。
对于某些计算，变化快慢反而更容易求。

\medskip
\textbf{\textcolor{CExplanation}{三、为什么不能随时分别求导？}}

\[
\lim_{x\to1}\frac{x}{x-1}
\]
不能代入，但它是“$\frac10$”，不是“$\frac00$”。
若生硬地分别求导，会得到 $\frac11=1$；而原分式从右边靠近 $1$ 时越来越大，
从左边靠近时越来越小，根本没有双侧的确定极限。
\textbf{所以：不能代入，绝不等于可以用洛必达法则。}

\medskip
\begin{minipage}{\linewidth}
\textbf{\textcolor{CExplanation}{四、考试中如何选择最快路线？}}

\begin{description}[style=nextline, leftmargin=2.5em, labelsep=0.4em, itemsep=0.35em]
    \item[\textbf{先直接代入：}] 若分母不为 $0$，通常直接算，不必求导。
    \item[\textbf{有共同因子：}] 像 $\frac{x^2-1}{x-1}$，先拆成 $(x-1)(x+1)$，约掉更快。
    \item[\textbf{含根号：}] 像 $\frac{\sqrt{x}-2}{x-4}$，先乘 $\sqrt{x}+2$，往往一步化简。
    \item[\textbf{拆不开、难化简：}] 检查是否为“$0$ 比 $0$”或“两边都无限增大”，并逐项确认求导条件，再考虑分别求导。
    \item[\textbf{求导一次还不能算：}] 若仍是允许的情况，可再求一次；若已不是允许的情况，立即停止，改用正常运算。
\end{description}
\end{minipage}

\medskip
\textbf{\textcolor{CExplanation}{五、两个容易忽略的发现}}

\textbf{发现一：同样趋于 $0$，结果不一定是有限数。}
例如 $\frac{x-1}{(x-1)^2}=\frac1{x-1}$，从右侧接近 $1$ 时会越来越大，
从左侧接近时会越来越小。左右方向可能给出不同结果。

\textbf{发现二：求导后的比值没有极限，不代表原式也没有。}
例如
\[
\frac{x+\sin x}{x}=1+\frac{\sin x}{x}\longrightarrow1
\qquad(x\to+\infty),
\]
因为 $-1\le\sin x\le1$；但分别求导后得到 $1+\cos x$，它一直变化，不能接近一个固定值。
此时只是\textbf{这条计算捷径走不通}，不是原题没有答案。



\textbf{\textcolor{CProof}{先说明证明的范围}}

洛必达法则可以处理比下面更复杂的情况，例如 $f'(a)=g'(a)=0$，甚至在 $a$ 处根本没有函数值。
这些情况的一般证明不能只凭高中导数定义直接推出。
这里不用超出高中范围的定理硬凑完整证明，而是用\textbf{三条可以亲手检查的路线}：
前两条证明两类常见的分式，第三条证明一个常用的三角函数极限。

\medskip

\end{explanationbox}
